\begingroup
\renewcommand{\thesection}{6B}
\section{Conservative Informative Action}\label{sec:conservativeaction}
\status{definition within the existing admissible-transduction framework}
Coupling specifies a dependence; action realizes a continuation; current additionally requires an encoded distinction, admitted field handoffs, and downstream decoding. Conservation is a property of the complete update, not a fourth name for any of these objects.

\begin{definition}[Conservative informative action]\label{def:conservativeaction}
Fix a complete registered state domain, its admitted recoverable transductions, a readout family that witnesses informativeness, and a declared invariant $I$. An informative action is conservative relative to $I$ when its realized transduction $\mathsf T$ satisfies $I(\mathsf Tz)=I(z)$ on its admitted domain.
\end{definition}
The invariant may be one readout of a much richer object. It need not equal $\mathcal C$, and invariance alone neither certifies recovery nor supplies an informative effect. In a faithful linear realization a useful special case is $I_{\mathsf G}(X)=X^\dagger\mathsf G X$. Section~\ref{sec:conservativedynamics} proves the associated skew-generator condition and its limits. ``Informational form'' is justified only when the representation and its relation to registered distinctions have also been declared.

Gyroscopic orientation and source-free electromagnetic evolution provide two conventional comparison models. They show how an invariant can constrain a changing state; they do not establish that every admissible QR action is linear, metric-preserving, or physical. Likewise an internal conservative update need not be an adjacent handoff and does not by itself increment propagation depth.
\endgroup
\setcounter{section}{6}
